% Section 10 - Annexes
\section{Annexes}

\subsection{Raccourcis clavier}

L'application propose plusieurs raccourcis clavier pour accélérer le travail.

\subsubsection{Raccourcis généraux}

\begin{longtable}{|p{4cm}|p{10cm}|}
\hline
\textbf{Raccourci} & \textbf{Action} \\
\hline
\endfirsthead
\hline
\textbf{Raccourci} & \textbf{Action} \\
\hline
\endhead
\hline
\endfoot

Ctrl + N & Créer une nouvelle configuration de classe \\
\hline
Ctrl + O & Ouvrir une configuration existante \\
\hline
Ctrl + S & Sauvegarder la configuration en cours \\
\hline
Ctrl + E & Charger un fichier Excel \\
\hline
Ctrl + Q & Quitter l'application \\
\hline
F5 & Actualiser l'interface \\
\hline
F1 & Afficher l'aide \\
\hline
F12 & Ouvrir la console de développement (avancé) \\
\hline
\end{longtable}

\subsubsection{Raccourcis spécifiques}

\begin{longtable}{|p{4cm}|p{10cm}|}
\hline
\textbf{Raccourci} & \textbf{Action} \\
\hline
\endfirsthead
\hline
\textbf{Raccourci} & \textbf{Action} \\
\hline
\endhead
\hline
\endfoot

Ctrl + G & Générer les relevés de notes \\
\hline
Ctrl + A & Générer les attestations \\
\hline
Ctrl + P & Prévisualiser un document \\
\hline
Ctrl + Shift + E & Export de la configuration \\
\hline
Ctrl + Shift + I & Import de configuration \\
\hline
\end{longtable}

\subsection{Format des codes QR}

\subsubsection{Contenu standard}

Les codes QR générés contiennent les informations au format texte brut :

\begin{verbatim}
MATRICULE: 2023001
NOM: DUPONT
PRENOM: Jean
SEMESTRE: S1
NIVEAU: L1
FILIERE: MEDECINE
\end{verbatim}

\subsubsection{Chiffrement}

Lorsque le chiffrement est activé :

\begin{itemize}[leftmargin=*]
    \item \textbf{Algorithme} : AES-128-ECB
    \item \textbf{Clé} : Basée sur le matricule de l'étudiant
    \item \textbf{Format} : Texte chiffré en Base64
    \item \textbf{Avantage} : Impossible de falsifier sans la clé
\end{itemize}

\subsubsection{Niveaux de correction d'erreur}

\begin{table}[h]
\centering
\begin{tabular}{|c|c|p{8cm}|}
\hline
\textbf{Niveau} & \textbf{Récupération} & \textbf{Usage recommandé} \\
\hline
L & 7\% & Documents de qualité, impression professionnelle \\
\hline
M & 15\% & Usage standard (recommandé) \\
\hline
Q & 25\% & Documents photocopiés, archivage long terme \\
\hline
H & 30\% & Environnements difficiles, manipulations fréquentes \\
\hline
\end{tabular}
\caption{Niveaux de correction d'erreur des QR codes}
\end{table}

\subsection{Glossaire}

\begin{description}[leftmargin=5cm, style=nextline]
    \item[AES] Advanced Encryption Standard - Algorithme de chiffrement symétrique

    \item[Attestation de réussite] Document officiel certifiant qu'un étudiant a validé un semestre (moyenne $\geq$ 10/20)

    \item[Base de notation] Échelle sur laquelle les notes sont données (10, 20, ou 100)

    \item[Chiffrement compact] Méthode de chiffrement produisant une sortie de taille réduite, adaptée aux QR codes

    \item[Coefficient] Poids d'un EC dans le calcul de la moyenne d'une UE

    \item[Configuration de classe] Ensemble structuré définissant les semestres, UEs et ECs d'une classe

    \item[Crédits ECTS] European Credit Transfer and Accumulation System - Système européen de crédits académiques

    \item[EC] Élément Constitutif - Matière ou cours au sein d'une UE

    \item[Excel] Format de fichier tableur (.xlsx, .xls) utilisé pour importer les données

    \item[Mapping] Association entre les colonnes Excel et les ECs de la configuration

    \item[Matricule] Numéro d'identification unique d'un étudiant

    \item[Mention] Appréciation qualitative basée sur la moyenne (Très Bien, Bien, Assez Bien, Passable, Ajourné)

    \item[Mode démo] Mode de fonctionnement sans licence, avec restrictions

    \item[Multi-semestres] Fonctionnalité permettant de générer des documents pour plusieurs semestres simultanément

    \item[OOM] Out Of Memory - Erreur survenant quand la mémoire RAM est insuffisante

    \item[PDF fusionné] Un seul fichier PDF contenant plusieurs documents assemblés séquentiellement

    \item[Préréglage] Thème prédéfini prêt à l'emploi

    \item[QR code] Quick Response Code - Code-barres 2D lisible par smartphone

    \item[Relevé de notes] Document officiel récapitulant les notes, moyennes et crédits d'un étudiant pour un semestre

    \item[Semestre] Période académique (généralement 6 mois) regroupant plusieurs UEs

    \item[Thème] Ensemble de paramètres visuels (couleurs, polices, mise en page) appliqué aux documents

    \item[UE] Unité d'Enseignement - Regroupement cohérent d'ECs valant un certain nombre de crédits

    \item[ZIP] Format d'archive compressée permettant de regrouper plusieurs fichiers
\end{description}

\subsection{Limites techniques}

\subsubsection{Limites pour les relevés de notes}

\begin{itemize}[leftmargin=*]
    \item Nombre maximum de relevés par génération : \textbf{500} (recommandé : 200)
    \item Nombre maximum de semestres par configuration : \textbf{12}
    \item Nombre maximum d'UEs par semestre : \textbf{20}
    \item Nombre maximum d'ECs par UE : \textbf{15}
    \item Taille maximale du fichier Excel : \textbf{50 Mo}
    \item Nombre maximum de lignes Excel : \textbf{10 000}
\end{itemize}

\subsubsection{Limites pour les attestations}

\begin{itemize}[leftmargin=*]
    \item Nombre maximum d'attestations par génération : \textbf{1000} (recommandé : 500)
    \item Nombre maximum de thèmes personnalisés sauvegardés : \textbf{50}
    \item Taille maximale du logo : \textbf{2 Mo} (format PNG ou JPG)
    \item Résolution logo recommandée : \textbf{800x200 pixels}
\end{itemize}

\subsubsection{Limites pour les QR codes}

\begin{itemize}[leftmargin=*]
    \item Nombre maximum de PDFs traités simultanément : \textbf{200}
    \item Taille maximale d'un PDF individuel : \textbf{50 Mo}
    \item Taille QR code : Min \textbf{40 px}, Max \textbf{200 px}
    \item Nombre de colonnes dans QR : Recommandé max \textbf{6}
    \item Longueur totale contenu QR : Max \textbf{2000 caractères}
\end{itemize}

\subsection{Compatibilité des formats}

\subsubsection{Formats de fichiers en entrée}

\begin{table}[h]
\centering
\begin{tabular}{|l|p{5cm}|p{5cm}|}
\hline
\textbf{Type} & \textbf{Formats acceptés} & \textbf{Notes} \\
\hline
Données étudiants & .xlsx, .xls, .csv & Excel 2007+ recommandé \\
\hline
Documents PDF & .pdf & Version PDF 1.4 à 1.7 \\
\hline
Images (logo) & .png, .jpg, .jpeg & PNG avec transparence recommandé \\
\hline
Configurations & .json & Export de l'application \\
\hline
Thèmes & .json & Export de l'application \\
\hline
\end{tabular}
\caption{Formats de fichiers supportés en entrée}
\end{table}

\subsubsection{Formats de fichiers en sortie}

\begin{table}[h]
\centering
\begin{tabular}{|l|p{10cm}|}
\hline
\textbf{Format} & \textbf{Description} \\
\hline
PDF & Documents générés (relevés, attestations) au format PDF/A \\
\hline
ZIP & Archives compressées (compression DEFLATE) \\
\hline
JSON & Exports de configurations, thèmes, historique \\
\hline
\end{tabular}
\caption{Formats de fichiers générés}
\end{table}

\subsection{FAQ - Foire aux questions}

\subsubsection{Questions générales}

\textbf{Q : L'application nécessite-t-elle une connexion Internet ?}

R : Non. L'application fonctionne entièrement hors ligne. Seule la vérification des mises à jour nécessite Internet (optionnel).

\textbf{Q : Mes données sont-elles sécurisées ?}

R : Oui. Toutes les données restent sur votre ordinateur local. Rien n'est envoyé sur Internet.

\textbf{Q : Puis-je utiliser l'application sur plusieurs ordinateurs ?}

R : Oui, mais vous devez installer l'application sur chaque ordinateur. Les configurations peuvent être partagées via export/import.

\textbf{Q : L'application est-elle compatible avec macOS et Linux ?}

R : Oui. Des versions sont disponibles pour Windows, macOS et Linux.

\subsubsection{Questions sur les relevés}

\textbf{Q : Comment ajouter une nouvelle UE à une configuration existante ?}

R : Menu \menu{Config des relevés} > Sélectionnez votre configuration > Cliquez sur le semestre > \bouton{Ajouter une UE}.

\textbf{Q : Les moyennes sont arrondies comment ?}

R : Deux décimales. Exemple : 12.456 devient 12.46.

\textbf{Q : Que se passe-t-il si une note est manquante ?}

R : La note est ignorée dans le calcul. Seules les notes présentes sont prises en compte.

\textbf{Q : Puis-je changer le barème des mentions ?}

R : Non, le barème est fixe : Très Bien ($\geq$16), Bien (14-15.99), Assez Bien (12-13.99), Passable (10-11.99), Ajourné (<10).

\subsubsection{Questions sur les attestations}

\textbf{Q : Pourquoi un étudiant avec moyenne 9.9 ne reçoit pas d'attestation ?}

R : Les attestations de réussite sont réservées aux étudiants admis (moyenne $\geq$ 10/20). C'est un filtrage automatique.

\textbf{Q : Puis-je forcer la génération d'une attestation pour un étudiant ajourné ?}

R : Non. L'application n'autorise que les étudiants ayant validé leur semestre.

\textbf{Q : Combien de thèmes puis-je créer ?}

R : Maximum 50 thèmes personnalisés sauvegardés.

\textbf{Q : Les attestations bilingues nécessitent-elles deux générations ?}

R : Non. Une seule génération. Sélectionnez la langue dans les paramètres, l'application utilise les bonnes colonnes Excel.

\subsubsection{Questions sur les QR codes}

\textbf{Q : Le QR code est-il ajouté sur toutes les pages ou seulement la première ?}

R : Sur \textbf{TOUTES les pages} de chaque document. C'est une caractéristique unique de cette application.

\textbf{Q : Puis-je scanner un QR code directement depuis l'écran ?}

R : Oui. Ouvrez le PDF sur votre ordinateur et scannez avec un smartphone.

\textbf{Q : Quelle application de scan recommandez-vous ?}

R : iOS : Appareil photo natif. Android : Google Lens. Ou toute app "QR Code Reader".

\textbf{Q : Le QR code chiffré peut-il être déchiffré ?}

R : Uniquement avec la clé appropriée (basée sur le matricule de l'étudiant).

\subsection{Support et ressources}

\subsubsection{Canaux de support}

\begin{itemize}[leftmargin=*]
    \item \textbf{Email} : support@fmsp.edu
    \item \textbf{Documentation en ligne} : https://docs.fmsp.edu/releves
    \item \textbf{Système de tickets} : https://support.fmsp.edu
    \item \textbf{Menu Aide} : Intégré dans l'application (F1)
\end{itemize}

\subsubsection{Ressources utiles}

\begin{itemize}[leftmargin=*]
    \item Guide de démarrage rapide (PDF) : Disponible dans le menu \menu{Aide}
    \item Tutoriels vidéo : https://videos.fmsp.edu/tutoriels
    \item Templates Excel : https://download.fmsp.edu/templates
    \item Forum communautaire : https://forum.fmsp.edu
\end{itemize}

\subsubsection{Signaler un bug}

Pour signaler un problème technique :

\begin{enumerate}[leftmargin=*]
    \item Menu \menu{Aide} > \bouton{Rapporter un problème}
    \item Ou envoyez un email à bugs@fmsp.edu
    \item Incluez :
    \begin{itemize}
        \item Version de l'application
        \item Description détaillée
        \item Étapes pour reproduire
        \item Captures d'écran
        \item Fichiers de logs (si applicable)
    \end{itemize}
\end{enumerate}

\subsection{Historique des versions}

\subsubsection{Version 2.0 (Actuelle)}

\textbf{Nouvelles fonctionnalités :}
\begin{itemize}[leftmargin=*]
    \item Génération multi-semestres améliorée
    \item Support bilingue FR/EN pour attestations
    \item 6 nouveaux préréglages de thèmes
    \item QR codes sur toutes les pages des PDFs multi-pages
    \item Export de configurations et thèmes
    \item Historique des documents générés
\end{itemize}

\textbf{Améliorations :}
\begin{itemize}[leftmargin=*]
    \item Optimisation mémoire pour générations massives
    \item Interface utilisateur repensée
    \item Performances de génération PDF améliorées de 30\%
    \item Meilleure gestion des erreurs Excel
\end{itemize}

\subsubsection{Version 1.5}

\textbf{Fonctionnalités :}
\begin{itemize}[leftmargin=*]
    \item Première version avec QR codes chiffrés
    \item Thèmes personnalisables
    \item Export ZIP avec compression
\end{itemize}

\subsubsection{Version 1.0 (Initiale)}

\textbf{Fonctionnalités de base :}
\begin{itemize}[leftmargin=*]
    \item Génération de relevés de notes
    \item Génération d'attestations de réussite
    \item Configuration académique basique
    \item Export PDF simple
\end{itemize}

\vfill

\begin{center}
\rule{\textwidth}{0.4pt}

\vspace{0.5cm}

\textit{Fin du guide d'utilisation complet}

\vspace{0.3cm}

\textbf{Version 2.0 - Corrigée et détaillée}

\vspace{0.2cm}

Basée sur l'analyse complète du code source

\vspace{0.3cm}

\today

\vspace{0.5cm}

\textit{Faculté de Médecine et Sciences Pharmaceutiques}

\vspace{0.2cm}

Pour toute question : support@fmsp.edu

\vspace{0.5cm}

\rule{\textwidth}{0.4pt}
\end{center}

\end{document}
